% GENERATED FILE. Edit canonical lessons, then run npm run generate:books.
% canonical-source-hash: fnv1a64:226afd1f5d2f3a90
% canonical-lessons: BN-C461-jedi, BN-C461-ekguye, BN-C461-chonchol, BN-C461-khamkheyali, BN-C461-go-dhora, BN-R461-review-454-461

\chapter{Stubborn and Restless}
\label{ch:bn-stubborn-and-restless}
\hlchaptermodality{461}

\begin{chapteropening}
\textbf{By the end of this chapter:} \emph{I can say that someone is stubborn, restless or whimsical.}

Complete the last lesson of chapter 461: I can say that someone is stubborn, restless or whimsical.
\end{chapteropening}

\section[\texorpdfstring{\bn{জেদি} (jedi)}{জেদি (jedi)}]{\bn{জেদি} (jedi) — headstrong; stubborn}

\label{lesson:BN-C461-jedi}



\begin{quote}
Before the new one: say the Bengali for selfish, then the Bengali for to oppress.
\end{quote}

\subsection*{You'll want to know: \bn{জেদি}}
\textbf{\bn{জেদি}} — \emph{jedi} — \textquotedblleft{}headstrong; stubborn\textquotedblright{}.

\begin{quote}
\bn{আমার} \bn{জেদি} \bn{বোন} \bn{খেলনাটা} \bn{ছাড়ল} \bn{না।}
\end{quote}

\begin{quote}
\emph{āmār jedi bon khelnāṭā chhāṛlô nā.}
\end{quote}

\begin{quote}
My headstrong sister would not let go of the toy.
\end{quote}

\begin{cousinweb}
\textbf{\bn{জেদি}} (\emph{jedi}) is \emph{jed}, \textquotedblleft{}stubbornness\textquotedblright{}, usually traced to Arabic \emph{ḍidd}, \textquotedblleft{}opposition\textquotedblright{}, through Persian, with the ending \emph{-i}.
\end{cousinweb}

\begin{grammarlens}[title={What you've built}]
One more word for describing what you lived and felt.
\end{grammarlens}

\subsection*{Your turn}
\begin{itemize}
\raggedright
  \item \emph{Say it:} \emph{jedi}
  \item \emph{Say it:} \emph{jedi}, once more
  \item \emph{Say it:} \emph{āmār jedi bon khelnāṭā chhāṛlô nā.} — My headstrong sister would not let go of the toy.
  \item \emph{Recall:} say the Bengali for selfish, then the Bengali for to oppress, then say \emph{jedi} again
\end{itemize}

\begin{tcolorbox}[breakable,colback=teal!4,colframe=teal!35!black,title={Before you move on}]
What does \bn{জেদি} mean? (\textquotedblleft{}Headstrong; stubborn\textquotedblright{}.) Say it once more.
\end{tcolorbox}
\section[\texorpdfstring{\bn{একগুঁয়ে} (ekgũye)}{একগুঁয়ে (ekgũye)}]{\bn{একগুঁয়ে} (ekgũye) — obstinate; pig-headed}

\label{lesson:BN-C461-ekguye}



\begin{quote}
Before the new one: say the Bengali for to oppress, then the Bengali for headstrong.
\end{quote}

\subsection*{You'll want to know: \bn{একগুঁয়ে}}
\textbf{\bn{একগুঁয়ে}} — \emph{ekgũye} — \textquotedblleft{}obstinate; pig-headed\textquotedblright{}.

\begin{quote}
\bn{একগুঁয়ে} \bn{ছেলেটা} \bn{ডাক্তারের} \bn{কাছে} \bn{যেতে} \bn{রাজি} \bn{হল} \bn{না।}
\end{quote}

\begin{quote}
\emph{ekgũye chheleṭā ḍāktārer kāchhe jete rāji hôlô nā.}
\end{quote}

\begin{quote}
The obstinate boy would not agree to go to the doctor.
\end{quote}

\begin{cousinweb}
\textbf{\bn{একগুঁয়ে}} (\emph{ekgũye}) begins with \bn{এক}, one, as if set on one thing only; the origin of the second part is not settled.
\end{cousinweb}

\begin{grammarlens}[title={What you've built}]
One more word for describing what you lived and felt.
\end{grammarlens}

\subsection*{Your turn}
\begin{itemize}
\raggedright
  \item \emph{Say it:} \emph{ekgũye}
  \item \emph{Say it:} \emph{ekgũye}, once more
  \item \emph{Say it:} \emph{ekgũye chheleṭā ḍāktārer kāchhe jete rāji hôlô nā.} — The obstinate boy would not agree to go to the doctor.
  \item \emph{Recall:} say the Bengali for to oppress, then the Bengali for headstrong, then say \emph{ekgũye} again
\end{itemize}

\begin{tcolorbox}[breakable,colback=teal!4,colframe=teal!35!black,title={Before you move on}]
What does \bn{একগুঁয়ে} mean? (\textquotedblleft{}Obstinate; pig-headed\textquotedblright{}.) Say it once more.
\end{tcolorbox}
\section[\texorpdfstring{\bn{চঞ্চল} (chônchôl)}{চঞ্চল (chônchôl)}]{\bn{চঞ্চল} (chônchôl) — restless; never still}

\label{lesson:BN-C461-chonchol}



\begin{quote}
Before the new one: say the Bengali for headstrong, then the Bengali for obstinate.
\end{quote}

\subsection*{You'll want to know: \bn{চঞ্চল}}
\textbf{\bn{চঞ্চল}} — \emph{chônchôl} — \textquotedblleft{}restless; never still\textquotedblright{}.

\begin{quote}
\bn{চঞ্চল} \bn{শিশুটা} \bn{এক} \bn{জায়গায়} \bn{বসে} \bn{থাকে} \bn{না।}
\end{quote}

\begin{quote}
\emph{chônchôl shishuṭā ek jāygāy bôse thāke nā.}
\end{quote}

\begin{quote}
The fidgety child never stays sitting in one place.
\end{quote}

\begin{cousinweb}
\textbf{\bn{চঞ্চল}} (\emph{chônchôl}) is a tatsama, Sanskrit \emph{cañcala}, \textquotedblleft{}moving to and fro, unsteady, flickering\textquotedblright{} (Monier-Williams).
\end{cousinweb}

\begin{grammarlens}[title={What you've built}]
One more word for describing what you lived and felt.
\end{grammarlens}

\subsection*{Your turn}
\begin{itemize}
\raggedright
  \item \emph{Say it:} \emph{chônchôl}
  \item \emph{Say it:} \emph{chônchôl}, once more
  \item \emph{Say it:} \emph{chônchôl shishuṭā ek jāygāy bôse thāke nā.} — The fidgety child never stays sitting in one place.
  \item \emph{Recall:} say the Bengali for headstrong, then the Bengali for obstinate, then say \emph{chônchôl} again
\end{itemize}

\begin{tcolorbox}[breakable,colback=teal!4,colframe=teal!35!black,title={Before you move on}]
What does \bn{চঞ্চল} mean? (\textquotedblleft{}Restless; never still\textquotedblright{}.) Say it once more.
\end{tcolorbox}
\section[\texorpdfstring{\bn{খামখেয়ালি} (khāmkheyāli)}{খামখেয়ালি (khāmkheyāli)}]{\bn{খামখেয়ালি} (khāmkheyāli) — whimsical; capricious}

\label{lesson:BN-C461-khamkheyali}



\begin{quote}
Before the new one: say the Bengali for obstinate, then the Bengali for never still.
\end{quote}

\subsection*{You'll want to know: \bn{খামখেয়ালি}}
\textbf{\bn{খামখেয়ালি}} — \emph{khāmkheyāli} — \textquotedblleft{}whimsical; capricious\textquotedblright{}.

\begin{quote}
\bn{আমার} \bn{খামখেয়ালি} \bn{মামা} \bn{হঠাৎ} \bn{চাকরি} \bn{ছাড়লেন।}
\end{quote}

\begin{quote}
\emph{āmār khāmkheyāli māmā hôṭhāt chākri chhāṛlen.}
\end{quote}

\begin{quote}
My whimsical uncle suddenly quit his job.
\end{quote}

\begin{cousinweb}
\textbf{\bn{খামখেয়ালি}} (\emph{khāmkheyāli}) is usually traced to Persian \emph{khām}, \textquotedblleft{}raw, unripe\textquotedblright{}, and \emph{khayāl}, \textquotedblleft{}a fancy\textquotedblright{} (Arabic, through Persian), the \bn{খেয়াল} of \bn{খেয়াল} \bn{করা} (taught): a half-baked fancy.
\end{cousinweb}

\begin{grammarlens}[title={What you've built}]
One more word for describing what you lived and felt.
\end{grammarlens}

\subsection*{Your turn}
\begin{itemize}
\raggedright
  \item \emph{Say it:} \emph{khāmkheyāli}
  \item \emph{Say it:} \emph{khāmkheyāli}, once more
  \item \emph{Say it:} \emph{āmār khāmkheyāli māmā hôṭhāt chākri chhāṛlen.} — My whimsical uncle suddenly quit his job.
  \item \emph{Recall:} say the Bengali for obstinate, then the Bengali for never still, then say \emph{khāmkheyāli} again
\end{itemize}

\begin{tcolorbox}[breakable,colback=teal!4,colframe=teal!35!black,title={Before you move on}]
What does \bn{খামখেয়ালি} mean? (\textquotedblleft{}Whimsical; capricious\textquotedblright{}.) Say it once more.
\end{tcolorbox}
\section[\texorpdfstring{\bn{গোঁ} \bn{ধরা} (gõ dhôrā)}{গোঁ ধরা (gõ dhôrā)}]{\bn{গোঁ} \bn{ধরা} (gõ dhôrā) — to dig one's heels in}

\label{lesson:BN-C461-go-dhora}



\begin{quote}
Before the new one: say the Bengali for never still, then the Bengali for whimsical.
\end{quote}

\subsection*{You'll want to know: \bn{গোঁ} \bn{ধরা}}
\textbf{\bn{গোঁ} \bn{ধরা}} — \emph{gõ dhôrā} — \textquotedblleft{}to dig one's heels in\textquotedblright{}.

\begin{quote}
\bn{নতুন} \bn{সাইকেলের} \bn{জন্য} \bn{ছেলেটা} \bn{গোঁ} \bn{ধরল।}
\end{quote}

\begin{quote}
\emph{nôtun sāikeler jônyô chheleṭā gõ dhôrlô.}
\end{quote}

\begin{quote}
The boy dug his heels in about getting a new bicycle.
\end{quote}

Past (he/she …): \bn{তিনি} \bn{গোঁ} \bn{ধরলেন}, \emph{tini gõ dhôrlen}.

\begin{cousinweb}
\textbf{\bn{গোঁ} \bn{ধরা}} (\emph{gõ dhôrā}) is \emph{gõ}, a stubborn will, which CDIAL lists for Bengali as a whim (CDIAL 4025), with \bn{ধরা} (taught), to hold (CDIAL 6747).
\end{cousinweb}

\begin{grammarlens}[title={What you've built}]
One more word for describing what you lived and felt.
\end{grammarlens}

\subsection*{Your turn}
\begin{itemize}
\raggedright
  \item \emph{Say it:} \emph{gõ dhôrā}
  \item \emph{Say it:} \emph{gõ dhôrā}, once more
  \item \emph{Say it:} \emph{nôtun sāikeler jônyô chheleṭā gõ dhôrlô.} — The boy dug his heels in about getting a new bicycle.
  \item \emph{Recall:} say the Bengali for never still, then the Bengali for whimsical, then say \emph{gõ dhôrā} again
\end{itemize}

\begin{tcolorbox}[breakable,colback=teal!4,colframe=teal!35!black,title={Before you move on}]
What does \bn{গোঁ} \bn{ধরা} mean? (\textquotedblleft{}To dig one's heels in\textquotedblright{}.) Say it once more.
\end{tcolorbox}
\section[(dialogue)]{Looking back: from Shyness and Hesitation to Stubborn and Restless}

\label{lesson:BN-R461-review-454-461}



\begin{quote}
Say the Bengali for reluctance and being in two minds.
\end{quote}

\subsection*{Your turn}
Say these aloud:

\begin{itemize}
\raggedright
  \item reluctance, being in two minds, stiff and awkward, huddled up, to hesitate
  \item patience, enthusiasm, keen interest, eager, to be able to wait
  \item character, a good quality, a habit, conduct, to master
  \item generous, miserly, sincere, helpful to others, noble
  \item honest, two-faced, trustworthy, cunning, to cheat
  \item modest, soft-spoken, insolent, ill-mannered, to flatter
  \item cruel, stern, aggressive, selfish, to oppress
  \item headstrong, obstinate, never still, whimsical, to dig one's heels in
\end{itemize}

\begin{tcolorbox}[breakable,colback=teal!4,colframe=teal!35!black,title={Before you move on}]
Reluctance? (\textbf{\bn{সংকোচ}}, \emph{sôngkoch}.) Being in two minds? (\textbf{\bn{দ্বিধা}}, \emph{dvidhā}.) Stiff and awkward? (\textbf{\bn{আড়ষ্ট}}, \emph{āṛôshṭô}.) Huddled up? (\textbf{\bn{জড়সড়}}, \emph{jôṛôsôṛô}.) To hesitate? (\textbf{\bn{ইতস্তত} \bn{করা}}, \emph{itôstôtô kôrā}.) Patience? (\textbf{\bn{ধৈর্য}}, \emph{dhoirjô}.) Enthusiasm? (\textbf{\bn{উৎসাহ}}, \emph{utsāhô}.) Keen interest? (\textbf{\bn{আগ্রহ}}, \emph{āgrôhô}.) Eager? (\textbf{\bn{উৎসুক}}, \emph{utsuk}.) To be able to wait? (\textbf{\bn{তর} \bn{সওয়া}}, \emph{tôr sôoyā}.) Character? (\textbf{\bn{চরিত্র}}, \emph{chôritrô}.) A good quality? (\textbf{\bn{গুণ}}, \emph{gun}.) A habit? (\textbf{\bn{অভ্যাস}}, \emph{ôbhyās}.) Conduct? (\textbf{\bn{আচরণ}}, \emph{āchôrôn}.) To master? (\textbf{\bn{রপ্ত} \bn{করা}}, \emph{rôptô kôrā}.) Generous? (\textbf{\bn{উদার}}, \emph{udār}.) Miserly? (\textbf{\bn{কৃপণ}}, \emph{kripôn}.) Sincere? (\textbf{\bn{আন্তরিক}}, \emph{āntôrik}.) Helpful to others? (\textbf{\bn{পরোপকারী}}, \emph{pôropôkārī}.) Noble? (\textbf{\bn{মহৎ}}, \emph{môhôt}.) Honest? (\textbf{\bn{সৎ}}, \emph{sôt}.) Two-faced? (\textbf{\bn{কপট}}, \emph{kôpôṭ}.) Trustworthy? (\textbf{\bn{বিশ্বস্ত}}, \emph{bishvôstô}.) Cunning? (\textbf{\bn{ধূর্ত}}, \emph{dhūrtô}.) To cheat? (\textbf{\bn{ঠকানো}}, \emph{ṭhôkāno}.) Modest? (\textbf{\bn{বিনয়ী}}, \emph{binôyī}.) Soft-spoken? (\textbf{\bn{নম্র}}, \emph{nômrô}.) Insolent? (\textbf{\bn{উদ্ধত}}, \emph{uddhôtô}.) Ill-mannered? (\textbf{\bn{বেয়াদব}}, \emph{beyādôb}.) To flatter? (\textbf{\bn{তোষামোদ} \bn{করা}}, \emph{toshāmod kôrā}.) Cruel? (\textbf{\bn{নিষ্ঠুর}}, \emph{nishṭhur}.) Stern? (\textbf{\bn{কঠোর}}, \emph{kôṭhor}.) Aggressive? (\textbf{\bn{উগ্র}}, \emph{ugrô}.) Selfish? (\textbf{\bn{স্বার্থপর}}, \emph{svārthôpôr}.) To oppress? (\textbf{\bn{অত্যাচার} \bn{করা}}, \emph{ôtyāchār kôrā}.) Headstrong? (\textbf{\bn{জেদি}}, \emph{jedi}.) Obstinate? (\textbf{\bn{একগুঁয়ে}}, \emph{ekgũye}.) Never still? (\textbf{\bn{চঞ্চল}}, \emph{chônchôl}.) Whimsical? (\textbf{\bn{খামখেয়ালি}}, \emph{khāmkheyāli}.) To dig one's heels in? (\textbf{\bn{গোঁ} \bn{ধরা}}, \emph{gõ dhôrā}.)
\end{tcolorbox}
